%% ----------------------------------------------------------------------------
%%  §5.6  Significato geometrico della derivata. Retta tangente
%% ----------------------------------------------------------------------------
\section{Significato geometrico della derivata. Retta tangente}
\label{sec:retta-tangente}

Sia $f$ una funzione definita in un intorno di un punto $x_{0}$ e si consideri nel piano
$x$, $y$ il grafico della funzione, come nella \cref{fig:secanti}. Ci proponiamo di
determinare l'equazione della retta $r$ passante per il punto $P_{0}$ di coordinate
$(x_{0},f(x_{0}))$ e \emph{tangente} al grafico della funzione $f$. Per le circonferenze la
tangente si definisce come la retta che ha un solo punto in comune con la curva; per i
grafici di funzioni questa definizione non va bene (la retta $y=1$ ha infiniti punti in
comune con il grafico di $\sen x$, eppure è tangente in ciascuno di essi; l'asse $y$ ha un
solo punto in comune con il grafico di $x^{2}$, ma non è tangente), e occorre una
definizione basata sul concetto di limite.

Ciò che preliminarmente è più opportuno fare è determinare l'equazione di una retta $r'$
\emph{secante} il grafico della funzione $f$ nei punti $P_{0}=(x_{0},f(x_{0}))$ e
$P=(x_{0}+h,f(x_{0}+h))$, con $h\neq0$. L'equazione di una generica retta non verticale è
$y=mx+q$; determiniamo i parametri $m$, $q$ imponendo che la retta passi per i punti dati:
\begin{equation}\label{eq:sistema-secante}
  \begin{cases}
    f(x_{0})=mx_{0}+q & \text{(passaggio per } P_{0}),\\
    f(x_{0}+h)=m(x_{0}+h)+q & \text{(passaggio per } P).
  \end{cases}
\end{equation}
Abbiamo un sistema di due equazioni nelle due incognite $m$, $q$, che si può risolvere
sottraendo la prima equazione dalla seconda. Si ottiene $m=\frac{f(x_{0}+h)-f(x_{0})}{h}$,
e poi si ricava $q$ dalla prima equazione. L'equazione della retta secante risulta essere
\begin{equation}\label{eq:retta-secante}
  y=f(x_{0})+\frac{f(x_{0}+h)-f(x_{0})}{h}\,(x-x_{0}).
\end{equation}
Il coefficiente angolare della secante è dunque il rapporto incrementale di $f$ in $x_{0}$.

\begin{figure}[htb]
\centering
\begin{tikzpicture}
\begin{axis}[width=0.9\textwidth, height=6cm, xmin=0, xmax=6.3, ymin=0, ymax=5,
  axis lines=left, xtick={1.5}, xticklabels={$x_{0}$}, ytick={1.4}, yticklabels={$f(x_{0})$},
  xlabel={$x$}, ylabel={$y$}, clip=false]
  \addplot[curva, domain=0.5:6, samples=120] {1.4+1.2*(x-1.5)-0.16*(x-1.5)^2};
  \addplot[tratteggio] coordinates {(0,1.4) (1.5,1.4) (1.5,0)};
  \foreach \hh/\op/\xe in {4/0.35/6.1, 2.5/0.55/5.85, 1.2/0.8/4.95} {
    \edef\temp{\noexpand\addplot[curva2, opacity=\op, domain=0.6:\xe]
      {1.4+(1.2-0.16*\hh)*(x-1.5)};}
    \temp
    \edef\temp{\noexpand\addplot[only marks, mark=*, mark size=1.3pt, rame, opacity=\op]
      coordinates {(1.5+\hh, {1.4+1.2*\hh-0.16*\hh*\hh})};}
    \temp
  }
  \addplot[curva3, domain=0.6:4.4] {1.4+1.2*(x-1.5)};
  \addplot[only marks, mark=*, mark size=1.6pt, blu] coordinates {(1.5,1.4)};
  \node[font=\scriptsize, anchor=south east] at (axis cs:1.5,1.45) {$P_{0}$};
  \node[font=\scriptsize, verde!50!black, anchor=west] at (axis cs:4.4,4.88) {$r$};
  \node[font=\scriptsize, rame, anchor=west] at (axis cs:6.1,3.98) {$r'$};
\end{axis}
\end{tikzpicture}
\caption{Le rette secanti $r'$ per $P_{0}$ e per un secondo punto $P$ del grafico (in
rame, sempre più marcate al diminuire di $h$) tendono, quando $P$ si avvicina a $P_{0}$,
alla retta tangente $r$ (in verde), la cui pendenza è $f'(x_{0})$.}
\label{fig:secanti}
\end{figure}

L'equazione della retta tangente, quando esiste, è il limite per $h\to0$ dell'equazione
della retta secante. Si può passare al limite nella \eqref{eq:retta-secante} se e solo se
esiste finito il limite del rapporto incrementale, cioè se e solo se $f$ è derivabile in
$x_{0}$. Ciò motiva la definizione seguente.

\begin{definizione}[Retta tangente]\label{def:retta-tangente}\index{retta tangente}
Se $f$ è derivabile in $x_{0}$, si chiama \emph{retta tangente al grafico di $f$ nel
punto $(x_{0},f(x_{0}))$} la retta di equazione
\begin{equation}\label{eq:retta-tangente}
  y=f(x_{0})+f'(x_{0})(x-x_{0}).
\end{equation}
Se $f$ è continua in $x_{0}$ e il rapporto incrementale tende a $+\infty$ o a $-\infty$,
si dice che il grafico ha in $(x_{0},f(x_{0}))$ \emph{tangente verticale}, la retta
$x=x_{0}$.
\end{definizione}

Quanto stabilito fornisce il significato geometrico della derivata. Dato che
nell'equazione della retta tangente il coefficiente della $x$ è uguale a $m=f'(x_{0})$, si
dice che \emph{la derivata di una funzione $f$ in un punto $x_{0}$ è il coefficiente
angolare della retta tangente al grafico della funzione nel punto $(x_{0},f(x_{0}))$}. La
derivata è quindi una misura della \emph{pendenza} del grafico della funzione: se
$\theta\in\bigl(-\frac\pi2,\frac\pi2\bigr)$ è l'angolo che la tangente forma con l'asse $x$,
si ha $f'(x_{0})=\tg\theta$. Ad esempio, la tangente al grafico di $x^{2}$ nel punto $(1,1)$
ha coefficiente angolare $2$ ed equazione $y=1+2(x-1)=2x-1$.

\subsection{L'approssimazione lineare}

Diamo un esempio numerico di utilizzazione dell'equazione della retta tangente,
esaminando un problema di calcolo approssimato dei valori di una funzione. Normalmente non
è immediato il calcolo del valore numerico di una funzione in un punto. Ad esempio, è
facile calcolare a mente i valori numerici delle funzioni $\sqrt x$, oppure $\sen x$,
solo per particolari valori della $x$. Al contrario, è sempre elementare calcolare i
valori numerici delle funzioni $y=mx+q$, che hanno per grafico una retta. L'idea è quella
di «sostituire» una funzione data con la sua retta tangente in un punto di ascissa
$x_{0}$, con $x_{0}$ vicino al punto $x$ in cui si vuole calcolare la funzione. Dalla
\cref{fig:approssimazione-lineare} è intuitivamente chiaro che l'errore che si commette è
tanto più piccolo, quanto più $x$ è vicino all'ascissa del punto di tangenza $x_{0}$.
Cioè la quantità $f(x_{0})+f'(x_{0})(x-x_{0})$ rappresenta un'approssimazione di $f(x)$,
tanto migliore quanto più $x$ è vicino ad $x_{0}$; scriveremo
\begin{equation}\label{eq:approssimazione-lineare}
  f(x)\simeq f(x_{0})+f'(x_{0})(x-x_{0}) \qquad(\text{se } x \text{ è vicino ad } x_{0}).
\end{equation}
Il punto $x_{0}$ va scelto in modo che sia semplice calcolare $f(x_{0})$ e $f'(x_{0})$.

\begin{figure}[htb]
\centering
\begin{tikzpicture}
\begin{axis}[width=0.9\textwidth, height=5.6cm, xmin=0, xmax=6.2, ymin=0, ymax=4.3,
  axis lines=left, xtick={2.8,4.4}, xticklabels={$x_{0}$,$x$},
  ytick={2.7035,3.4023}, yticklabels={$t(x)$,$f(x)$},
  tick label style={font=\scriptsize}, xlabel={$x$}, ylabel={$y$}]
  \addplot[curva, domain=1:4.8, samples=100] {0.6+0.2*exp(0.6*x)};
  \addplot[curva3, domain=1:4.9] {1.6732+0.64387*(x-2.8)};
  \addplot[tratteggio] coordinates {(2.8,0) (2.8,1.6732)};
  \addplot[tratteggio] coordinates {(0,3.4023) (4.4,3.4023) (4.4,0)};
  \addplot[tratteggio] coordinates {(0,2.7035) (4.4,2.7035)};
  \draw[<->, rosso, thick] (axis cs:4.52,2.7035) -- node[pos=0.85, right, font=\scriptsize] {errore} (axis cs:4.52,3.4023);
  \addplot[only marks, mark=*, mark size=1.6pt, blu] coordinates {(2.8,1.6732)};
\end{axis}
\end{tikzpicture}
\caption{L'approssimazione lineare: vicino a $x_{0}$ il grafico di $f$ (blu) è ben
approssimato dalla retta tangente (verde), $t(x)=f(x_{0})+f'(x_{0})(x-x_{0})$. L'errore
$f(x)-t(x)$ tende a zero più rapidamente di $x-x_{0}$.}
\label{fig:approssimazione-lineare}
\end{figure}

Si può dare un significato rigoroso alla scrittura \eqref{eq:approssimazione-lineare}
usando i limiti. La \eqref{eq:approssimazione-lineare} significa che non solo la
differenza tra primo e secondo membro tende a zero quando $x\to x_{0}$ (questo varrebbe
per qualunque retta passante per $P_{0}$, se $f$ è continua), ma anche che tende a zero
\emph{più rapidamente} della quantità $x-x_{0}$. Conviene introdurre una notazione.

\begin{definizione}[Simbolo di Landau «o piccolo»]\label{def:o-piccolo}\index{o piccolo}
Siano $\phi$ e $\psi$ definite in un intorno di $x_{0}$ (escluso al più $x_{0}$), con
$\psi(x)\neq0$ per $x\neq x_{0}$. Si scrive $\phi(x)=o(\psi(x))$ per $x\to x_{0}$ (si
legge «$\phi$ è o piccolo di $\psi$») se
\[
  \lim_{x\to x_{0}}\frac{\phi(x)}{\psi(x)}=0 .
\]
\end{definizione}

Ad esempio $(x-x_{0})^{2}=o(x-x_{0})$ per $x\to x_{0}$, e $1-\cos x=o(x)$ per $x\to0$
\eqref{eq:lim-1-cos-x}.

\begin{proposizione}[Caratterizzazione della derivabilità]\label{prop:approssimazione-lineare}
Sia $f$ definita in un intorno di $x_{0}$. Allora $f$ è derivabile in $x_{0}$ se e solo
se esiste un numero $m\in\R$ tale che
\begin{equation}\label{eq:sviluppo-primo-ordine}
  f(x)=f(x_{0})+m(x-x_{0})+o(x-x_{0}) \qquad\text{per } x\to x_{0};
\end{equation}
in tal caso $m=f'(x_{0})$. In particolare, se $f$ è derivabile in $x_{0}$,
\begin{equation}\label{eq:errore-tangente}
  \lim_{x\to x_{0}}\frac{f(x)-[f(x_{0})+f'(x_{0})(x-x_{0})]}{x-x_{0}}=0 .
\end{equation}
\end{proposizione}

\begin{schema}
\Ipotesi{$f$ definita in un intorno di $x_{0}$.}
\Tesi{$f$ derivabile in $x_{0}$ $\sse$ esiste $m\in\R$ con
$f(x)=f(x_{0})+m(x-x_{0})+o(x-x_{0})$ per $x\to x_{0}$; \ in tal caso $m=f'(x_{0})$.}
\end{schema}

\begin{strategia}
Tutto si basa su un'identità algebrica: l'errore che si commette sostituendo $f$ con la
retta $y=f(x_{0})+m(x-x_{0})$, diviso per $x-x_{0}$, è uguale al rapporto incrementale
diminuito di $m$. Dire che l'errore è $o(x-x_{0})$ equivale quindi a dire che il rapporto
incrementale tende a $m$, cioè che $f$ è derivabile con $f'(x_{0})=m$: le due implicazioni
si leggono sulla stessa identità, nei due versi.
\end{strategia}

\begin{dimostrazione}
Per $x\neq x_{0}$ si ha l'identità
\begin{equation}\label{eq:identita-errore}
  \frac{f(x)-[f(x_{0})+m(x-x_{0})]}{x-x_{0}}=\frac{f(x)-f(x_{0})}{x-x_{0}}-m .
\end{equation}
Se $f$ è derivabile e $m=f'(x_{0})$, il secondo membro tende a $f'(x_{0})-f'(x_{0})=0$ per
la \eqref{eq:derivata-xi}: vale la \eqref{eq:sviluppo-primo-ordine}, e anche la
\eqref{eq:errore-tangente}. Viceversa, se vale la \eqref{eq:sviluppo-primo-ordine} per un
certo $m$, il primo membro della \eqref{eq:identita-errore} tende a zero, e quindi il
rapporto incrementale $\frac{f(x)-f(x_{0})}{x-x_{0}}$ tende a $m$: $f$ è derivabile e
$f'(x_{0})=m$.
\end{dimostrazione}

La proposizione dice che la retta tangente è \emph{l'unica} retta per $P_{0}$ che
approssima il grafico «al primo ordine», cioè con un errore che è $o(x-x_{0})$: qualunque
altra retta $y=f(x_{0})+m(x-x_{0})$ con $m\neq f'(x_{0})$ commette un errore il cui
rapporto con $x-x_{0}$ tende a $f'(x_{0})-m\neq0$. In questa forma la derivabilità si
estende alle funzioni di più variabili, come vedremo nel \cref{sec:differenziabilita}.
Si usa spesso anche la notazione del \emph{differenziale}: posto $dx=x-x_{0}$ e
$df=f'(x_{0})\,dx$, la \eqref{eq:sviluppo-primo-ordine} si scrive
$\Delta f=df+o(dx)$, e il differenziale $df$ è la parte lineare dell'incremento $\Delta f$.

Ad esempio, se $f(x)=\sqrt x$, la \eqref{eq:approssimazione-lineare} diventa
\begin{equation}\label{eq:approssimazione-radice}
  \sqrt x\simeq\sqrt{x_{0}}+\frac1{2\sqrt{x_{0}}}(x-x_{0}) \qquad(\text{se } x \text{ è vicino ad } x_{0}).
\end{equation}
Volendo esprimere in forma decimale $\sqrt{80}$, scegliendo $x_{0}=81$ otteniamo
\begin{equation}\label{eq:radice-80}
  \sqrt{80}\simeq9+\frac1{2\cdot9}(-1)=9-\frac1{18}=8{,}9444\ldots
\end{equation}
(il valore esatto è $\sqrt{80}=8{,}9442\ldots$). Qualcuno forse avrà pensato che il conto
è stato possibile soltanto perché $80$ è vicino al quadrato perfetto $81$. Proviamo con
$\sqrt2$: si può calcolare $\sqrt{200}$, e poi dividere il risultato per $10$. Il quadrato
più vicino a $200$ è $196=14^{2}$; si ha quindi
\begin{equation}\label{eq:radice-200}
  \sqrt{200}\simeq14+\frac1{2\cdot14}\cdot4=14+\frac17=14{,}1428\ldots
\end{equation}
Perciò $\sqrt2\simeq1{,}41428\ldots$ (il valore esatto di $\sqrt2$ è $1{,}41421\ldots$). I
conti fatti dovrebbero aver dato un'idea dell'utilità delle derivate nella tabulazione
delle funzioni elementari. Il programma che segue confronta l'approssimazione con il
valore vero in alcuni casi, e mostra che l'errore è, con buona approssimazione,
proporzionale a $(x-x_{0})^{2}$.

\begin{python}[Quanto è buona l'approssimazione lineare di $\sqrt x$?]
import math

# tangente a sqrt in x0: sqrt(x) ~ sqrt(x0) + (x - x0) / (2 sqrt(x0))
lin = lambda x, x0: math.sqrt(x0) + (x - x0) / (2 * math.sqrt(x0))
casi = [(80, 81), (200, 196), (102, 100), (110, 100), (150, 100)]
for x, x0 in casi:
    e = lin(x, x0) - math.sqrt(x)
    print(f"sqrt({x}) ~ {lin(x, x0):.6f}  vero {math.sqrt(x):.6f}  "
          f"err {e:.1e}  rapporto {e / (x - x0)**2:.5f}")
\end{python}
\begin{risultato}
sqrt(80) ~ 8.944444  vero 8.944272  err 1.7e-04  rapporto 0.00017
sqrt(200) ~ 14.142857  vero 14.142136  err 7.2e-04  rapporto 0.00005
sqrt(102) ~ 10.100000  vero 10.099505  err 5.0e-04  rapporto 0.00012
sqrt(110) ~ 10.500000  vero 10.488088  err 1.2e-02  rapporto 0.00012
sqrt(150) ~ 12.500000  vero 12.247449  err 2.5e-01  rapporto 0.00010
\end{risultato}

Con $x_{0}=100$, passando da $x=102$ a $x=110$ l'incremento $x-x_{0}$ diventa $5$ volte
più grande e l'errore circa $25$ volte più grande: è il segno che l'errore è dell'ordine
di $(x-x_{0})^{2}$, e il rapporto $\frac{\text{errore}}{(x-x_{0})^{2}}$ è quasi costante
(per $x_{0}=100$ vale circa $0{,}000125=\frac1{8\cdot100^{3/2}}$). Torneremo su questo
interessante aspetto del calcolo differenziale con la formula di Taylor, che spiegherà
il ruolo della derivata seconda in questa costante.

\begin{idea}[Sensibilità e propagazione degli errori]
L'approssimazione lineare è lo strumento con cui l'ingegnere valuta come un errore sui
dati si propaga ai risultati: se $y=f(x)$ e il dato $x$ è affetto da un piccolo errore
$\Delta x$, l'errore sul risultato è $\Delta y\simeq f'(x)\,\Delta x$. Riprendiamo
l'anemometro del \cref{cap:funzioni-continue}: la velocità si ricava dalla differenza di
pressione misurata, $V=\sqrt{2\Delta p/\rho}$, e per la \eqref{eq:derivata-radice} e la
regola della catena
\[
  \frac{dV}{d(\Delta p)}=\frac1{\sqrt{2\rho\,\Delta p}}=\frac1{\rho V}.
\]
Un errore di $5$~Pa nella misura della differenza di pressione, al livello del mare,
produce un errore di circa $5/(1{,}225\cdot20)\approx0{,}20$~m/s a $20$~m/s, ma soltanto
di circa $0{,}04$~m/s a $100$~m/s: la derivata quantifica la sensibilità che nel capitolo
precedente la sola continuità non sapeva misurare. Lo stesso principio sta alla base
della \emph{linearizzazione} in meccanica del volo: per piccole variazioni
dell'incidenza attorno a un valore di riferimento $\alpha_{0}$ si scrive
$C_{L}\simeq C_{L}(\alpha_{0})+C_{L\alpha}\,(\alpha-\alpha_{0})$, cioè si sostituisce la curva
di portanza con la sua retta tangente, e $C_{L\alpha}=\frac{dC_{L}}{d\alpha}(\alpha_{0})$ è
proprio una derivata.
\end{idea}

\subsection{Le funzioni $x^{n}\sen\frac1x$}

Utilizzando fra l'altro il significato geometrico della derivata, studiamo la continuità
e la derivabilità, nel punto $x_{0}=0$, della funzione $f_{n}$ ($n=0,1,2$) definita su $\R$
dalla formula
\begin{equation}\label{eq:fn}
  f_{n}(0)=0, \qquad f_{n}(x)=x^{n}\sen\frac1x \quad\text{se } x\neq0 .
\end{equation}
Per $n=0$ si ottiene la funzione che vale $\sen\frac1x$ per $x\neq0$ e $0$ per $x=0$; per
$n=1$ e $n=2$ le funzioni $x\sen\frac1x$ e $x^{2}\sen\frac1x$, prolungate con il valore
$0$ in $x=0$.

\begin{itemize}
  \item La funzione $f_{0}$ non è continua nel punto $x_{0}=0$, ma presenta una
  discontinuità di seconda specie: come risulta dalla \eqref{eq:sen-1-x}, non esiste il
  limite per $x\to0$ di $f_{0}(x)$. Inoltre $f_{0}$, non essendo continua in $x_{0}=0$, non
  è neanche derivabile in tale punto (\cref{teo:derivabile-continua}).
  \item La funzione $f_{1}$, prodotto del fattore infinitesimo $x$ per il fattore limitato
  $\sen\frac1x$, converge a zero per $x\to0$ (\cref{es:x-sen-1-x}); essendo $f_{1}(0)=0$,
  la funzione è continua (anche) in $x_{0}=0$. Però non risulta derivabile in tale punto,
  perché non esiste il limite del rapporto incrementale:
  \begin{equation}\label{eq:f1-rapporto}
    \lim_{h\to0}\frac{f_{1}(h)-f_{1}(0)}{h}=\lim_{h\to0}\frac{h\sen\frac1h}{h}=\lim_{h\to0}\sen\frac1h,
  \end{equation}
  che non esiste per la \eqref{eq:sen-1-x}.
  \item La funzione $f_{2}$ è derivabile (e quindi anche continua) anche per $x_{0}=0$, e
  la derivata vale $f_{2}'(0)=0$; infatti, per l'\cref{es:x-sen-1-x},
  \begin{equation}\label{eq:f2-derivata}
    f_{2}'(0)=\lim_{h\to0}\frac{f_{2}(h)-f_{2}(0)}{h}=\lim_{h\to0}h\sen\frac1h=0 .
  \end{equation}
\end{itemize}
I risultati trovati sono riassunti nella seguente tabella.
\begin{center}
\renewcommand{\arraystretch}{1.3}
\begin{tabular}{c c c}
\toprule
funzione & continua in $x_{0}=0$ & derivabile in $x_{0}=0$\\
\midrule
$f_{0}(x)$ & no & no\\
$f_{1}(x)$ & sì & no\\
$f_{2}(x)$ & sì & sì\\
\bottomrule
\end{tabular}
\end{center}

\begin{figure}[htb]
\centering
\includegraphics[width=\textwidth]{xn_sen}
\caption{I grafici di $f_{1}$ e di $f_{2}$ vicino all'origine, compresi rispettivamente tra
$y=\pm\abs x$ e tra $y=\pm x^{2}$; a destra la derivata $f_{2}'(x)$ per $x\neq0$, che oscilla
tra valori vicini a $-1$ e a $1$, mentre $f_{2}'(0)=0$ (punto blu)
(\texttt{python/fig\_xn\_sen.py}).}
\label{fig:xn-sen}
\end{figure}

Nella \cref{fig:xn-sen} sono rappresentati i grafici delle funzioni $f_{1}$ e $f_{2}$; la
continuità di $f_{1}$, $f_{2}$ nel punto $x_{0}=0$ corrisponde a grafici «vicini»
all'origine delle coordinate quando l'ascissa $x$ è «vicina» a $x_{0}=0$ (si noti che tale
proprietà non è verificata per il grafico della funzione $f_{0}$ nella
\cref{fig:sen-1-x}). Invece la derivabilità di $f_{2}$ e la non derivabilità di $f_{1}$
corrispondono al fatto che il grafico di $f_{2}$ ammette retta tangente anche nel punto
$x_{0}=0$ (dato che $f_{2}(0)=f_{2}'(0)=0$, l'equazione \eqref{eq:retta-tangente} della
retta tangente è $y=0$), mentre il grafico di $f_{1}$ non ha retta tangente nell'origine:
le rette secanti per l'origine e per il punto $(h,f_{1}(h))$ hanno coefficiente angolare
$\sen\frac1h$, che oscilla tra $-1$ e $1$ senza stabilizzarsi, e i grafici di $y=x$,
$y=-x$, con coefficienti angolari $+1$ e $-1$, danno un'idea dell'oscillazione. Tali
proprietà sono evidenziate dalle limitazioni
\begin{equation}\label{eq:f1-f2-limitazioni}
  -\abs x\le f_{1}(x)\le\abs x, \qquad -x^{2}\le f_{2}(x)\le x^{2} \qquad\perogni x\in\R .
\end{equation}

\begin{osservazione}[Una derivata che esiste ovunque ma non è continua]\label{oss:f2-non-C1}
La funzione $f_{2}$ è derivabile su tutto $\R$: per $x\neq0$, con le regole del prodotto e
della catena,
\[
  f_{2}'(x)=2x\sen\frac1x+x^{2}\cos\frac1x\cdot\Bigl(-\frac1{x^{2}}\Bigr)
  =2x\sen\frac1x-\cos\frac1x,
\]
e in $0$ vale la \eqref{eq:f2-derivata}. Ma la funzione derivata $f_{2}'$ \emph{non è
continua} in $0$: il primo addendo tende a $0$, il secondo, $\cos\frac1x$, non ha limite
per $x\to0$ (si ragiona come per $\sen\frac1x$), e quindi $f_{2}'(x)$ non ha limite per
$x\to0$ (parte destra della \cref{fig:xn-sen}). Dunque $f_{2}$ è derivabile, ma non è di
classe $C^{1}$: l'esistenza della derivata in ogni punto non garantisce la regolarità
della funzione derivata. In questo senso la $f_{2}$ è un esempio da tenere a mente ogni
volta che si è tentati di dare per scontato che una derivata sia continua.
\end{osservazione}
